\newcommand{\ExperimentPlatform}{Talapas, Slurm job 48205110 on n0187}
\newcommand{\ExperimentCPU}{AMD EPYC 7713 64-Core Processor}
\newcommand{\ExperimentCompiler}{GCC 8.5.0}
\newcommand{\RunStatus}{Measured on the compute node inside the recorded Slurm allocation. Allocation: 64 CPUs and 4 GiB RAM.}
\newcommand{\MaxSize}{33,554,432}
\newcommand{\TrialCount}{7}
\newcommand{\ConfigurationCount}{60}
\newcommand{\MeasurementCount}{420}
\newcommand{\ThreadCounts}{1, 2, 4, 8, 16, 32, 64}
\newcommand{\SerialTime}{15.61}
\newcommand{\TreeTime}{18.34}
\newcommand{\TreeSpeedup}{0.85}
\newcommand{\TreeThreads}{64}
\newcommand{\DoublingTime}{98.12}
\newcommand{\DoublingSpeedup}{0.16}
\newcommand{\DoublingThreads}{64}
\newcommand{\ProblemSizes}{2^{16}, 2^{20}, 2^{24}, 2^{25}}
\newcommand{\HighThreadFinding}{At 64 threads, tree and doubling speedups were 0.85 and 0.16, respectively.}
\newcommand{\FindingSummary}{Neither parallel method's median time beat serial at any tested input size. Lower total work alone did not guarantee proportional thread scaling.}
\newcommand{\MachineLimit}{The node has 128 physical cores across two sockets and eight NUMA domains. The shared first-touch policy and core binding can affect serial/parallel ratios; this experiment does not isolate placement effects or characterize other nodes.}
